\section{Conclusion}
This work introduced PGA-UNet, a compact CNN that incorporates a box prompt into encoder features and decoder skip attention. On the fixed BTXRD and FracAtlas test splits, the model achieved covering-prompt Dice scores of 0.782 and 0.727 at $512\times512$, with higher observed Dice than the tested conventional prompt baselines. The existing binary-prompt comparison and component ablations provide evidence for the evaluated conditioning design, with effects that depend on the dataset. PGA-UNet contains approximately 2.96 million parameters, and the reported timing results characterize its computational cost in the tested environment.

These findings are limited to simulated lesion-covering prompts, image-level evaluation, and the reported training configurations. Evaluation with clinician-drawn prompts remains future work. A training-seed stability analysis for PGA-UNet (Section~\ref{sec:seedstability}) reports two of four planned training seeds on the original fixed split, with the remaining two still training, and repeated comparative training runs for the evaluated baselines remain future work. The exploratory MedSigLIP gate (Section~\ref{sec:selfassess-results}) shows that a lightweight post-hoc filter can reduce box-level false positives on lesion-free images while keeping most positive boxes, without resolving image-level false-positive detection; developing a purpose-built, better-calibrated lightweight gate, together with a post-gate segmentation-quality evaluation, is left for future work.
